% =============================================================================
% 016 / 068
% Status: raw
% =============================================================================
% Notes:
%
% Source question:
% 有人覺得蘇錫常和皖南吳語區的散沙化和對母語的不熟悉比浙江嚴重很多，尤其是在青少年一代中，您贊同這個觀點嗎
% =============================================================================

\chapter[有人覺得蘇錫常和皖南吳語區的散沙化和對母語的不熟悉比浙江嚴重很多，尤其是在青少年一代中，您贊同這個觀點嗎]{有人覺得蘇錫常和皖南吳語區的散沙化和對母語的不熟悉比浙江嚴重很多，尤其是在青少年一代中，您贊同這個觀點嗎}
\qaanswer{
我不贊同這樣的說法。因為這些都無法量化並被統計出來。無論是蘇南、徽州還是浙江都面臨普通話的強力影響與世代置換。而且在我們這代人可見的未來，吳語徹底滅亡。討論誰比誰對母語更不熟悉跟討論用幾種方法殺死將要被烹飪的章魚一樣，對與受體而言意義不大。我倒是對「如何讓吳語延續下去」，「章魚如何逃脫廚師的魔爪回到大海」這類的議題更感興趣。
}

